%% generato automaticamente da tools/build_tex.py --- non modificare
\chapter{Determinanti di un sudoku}
\label{cap:sudoku}
\quaderno{Sudoku.dfw}{26 Apr 2009}{42}

Aprile 2009, il quaderno più antico insieme a \dfw{teoria.dfw}. La domanda è semplice e un po' oziosa, del genere che viene in mente guardando una griglia riempita: che cosa succede se la si tratta come una matrice? Quarantadue espressioni dopo, la risposta è che non succede granché — ma il percorso per stabilirlo è ordinato, e un fatto interessante c'è.

La griglia in esame è un sudoku valido: ogni riga, ogni colonna e ognuno dei nove blocchi contengono le cifre da $1$ a $9$ una volta sola. La verifica è stata rifatta durante il recupero e conferma.

\section{Il determinante della griglia intera}

Si comincia dalla domanda più diretta.

\begin{espr}{1}{User}{}
\[
\begin{vmatrix} 9 & 1 & 3 & 7 & 8 & 6 & 5 & 4 & 2 \\ 4 & 8 & 7 & 5 & 2 & 3 & 1 & 9 & 6 \\ 2 & 6 & 5 & 4 & 1 & 9 & 3 & 8 & 7 \\ 8 & 5 & 6 & 3 & 7 & 4 & 2 & 1 & 9 \\ 7 & 9 & 2 & 1 & 6 & 5 & 4 & 3 & 8 \\ 3 & 4 & 1 & 2 & 9 & 8 & 7 & 6 & 5 \\ 6 & 3 & 4 & 9 & 5 & 7 & 8 & 2 & 1 \\ 5 & 2 & 9 & 8 & 4 & 1 & 6 & 7 & 3 \\ 1 & 7 & 8 & 6 & 3 & 2 & 9 & 5 & 4 \end{vmatrix}
\]
\end{espr}
La griglia intera come matrice $9\times9$.

\begin{espr}{2}{Simp(\#1)}{}
\[
-2975130
\]
\end{espr}
Il suo determinante. È divisibile per $45$ — proprietà comune a tutti i sudoku — e in questo caso anche per $405$.


Il valore è $-2\,975\,130$, e ha una proprietà che vale la pena di far notare perché il quaderno non la nota: è divisibile per $45$, e la dimostrazione è di una riga. Poiché ogni riga della griglia somma a $45$, sostituendo alla prima colonna la somma di tutte le colonne si ottiene una colonna con tutti gli elementi uguali a $45$, che si raccoglie a fattore comune. Il determinante di un sudoku è dunque sempre un multiplo di $45$. In questo caso lo è anche di $405$: $-2\,975\,130=405\cdot(-7346)$, e la fattorizzazione completa è $-2\cdot3^4\cdot5\cdot3673$.

\section{I nove blocchi}

Il passo successivo è prendere i nove blocchi $3\times3$ — quelli che nel gioco devono contenere tutte le cifre — e calcolarne i determinanti uno per uno.

\begin{espr}{3}{User}{}
\[
\begin{vmatrix} 9 & 1 & 3 \\ 4 & 8 & 7 \\ 2 & 6 & 5 \end{vmatrix}
\]
\end{espr}
Il primo blocco: righe $1$--$3$, colonne $1$--$3$.

\begin{espr}{4}{User}{}
\[
0
\]
\end{espr}
Il determinante è nullo: le tre righe del blocco sono linearmente dipendenti.

\begin{espr}{5}{User}{}
\[
\begin{vmatrix} 7 & 8 & 6 \\ 5 & 2 & 3 \\ 4 & 1 & 9 \end{vmatrix}
\]
\end{espr}
Il secondo blocco.

\begin{espr}{6}{Simp(\#5)}{}
\[
-177
\]
\end{espr}
$-177$.

\begin{espr}{7}{User}{}
\[
\begin{vmatrix} 5 & 4 & 2 \\ 1 & 9 & 6 \\ 3 & 8 & 7 \end{vmatrix}
\]
\end{espr}
Il terzo blocco.

\begin{espr}{8}{Simp(\#7)}{}
\[
81
\]
\end{espr}
$81$.

\begin{espr}{9}{User}{}
\[
\begin{vmatrix} 8 & 5 & 6 \\ 7 & 9 & 2 \\ 3 & 4 & 1 \end{vmatrix}
\]
\end{espr}
Il quarto blocco.

\begin{espr}{10}{Simp(\#9)}{}
\[
9
\]
\end{espr}
$9$.

\begin{espr}{11}{User}{}
\[
\begin{vmatrix} 3 & 7 & 4 \\ 1 & 6 & 5 \\ 2 & 9 & 8 \end{vmatrix}
\]
\end{espr}
Il quinto blocco, quello centrale.

\begin{espr}{12}{User}{}
\[
11
\]
\end{espr}
$11$.

\begin{espr}{13}{User}{}
\[
\begin{vmatrix} 2 & 1 & 9 \\ 4 & 3 & 8 \\ 7 & 6 & 5 \end{vmatrix}
\]
\end{espr}
Il sesto blocco.

\begin{espr}{14}{Simp(\#13)}{}
\[
-3
\]
\end{espr}
$-3$.

\begin{espr}{15}{User}{}
\[
\begin{vmatrix} 6 & 3 & 4 \\ 5 & 2 & 9 \\ 1 & 7 & 8 \end{vmatrix}
\]
\end{espr}
Il settimo blocco.

\begin{espr}{16}{Simp(\#15)}{}
\[
-243
\]
\end{espr}
$-243$.

\begin{espr}{17}{User}{}
\[
\begin{vmatrix} 9 & 5 & 7 \\ 8 & 4 & 1 \\ 6 & 3 & 2 \end{vmatrix}
\]
\end{espr}
L'ottavo blocco.

\begin{espr}{18}{Simp(\#17)}{}
\[
-5
\]
\end{espr}
$-5$.

\begin{espr}{19}{User}{}
\[
\begin{vmatrix} 8 & 2 & 1 \\ 6 & 7 & 3 \\ 9 & 5 & 4 \end{vmatrix}
\]
\end{espr}
Il nono blocco.

\begin{espr}{20}{Simp(User)}{}
\[
77
\]
\end{espr}
$77$.

\begin{espr}{21}{User}{}
\[
0 + -177 + 81 + 9 + 11 + -3 + -243 + -5 + 77
\]
\end{espr}
La somma dei nove valori, scritta per esteso.

\begin{espr}{22}{Simp(\#21)}{}
\[
-250
\]
\end{espr}
Vale $-250$.

\begin{figure}[htbp]
  \centering
  % La griglia studiata in Sudoku.dfw e i determinanti dei nove blocchi.
\begin{tikzpicture}[scale=0.56, every node/.style={font=\small}]

  % ---- la griglia ---------------------------------------------------
  \begin{scope}
    \def\griglia{
      {9,1,3,7,8,6,5,4,2},
      {4,8,7,5,2,3,1,9,6},
      {2,6,5,4,1,9,3,8,7},
      {8,5,6,3,7,4,2,1,9},
      {7,9,2,1,6,5,4,3,8},
      {3,4,1,2,9,8,7,6,5},
      {6,3,4,9,5,7,8,2,1},
      {5,2,9,8,4,1,6,7,3},
      {1,7,8,6,3,2,9,5,4}}

    \foreach \bi in {0,1,2}{
      \foreach \bj in {0,1,2}{
        \pgfmathtruncatemacro{\tinta}{mod(\bi+\bj,2)}
        \ifnum\tinta=0
          \fill[grigiocarta] (3*\bj,-3*\bi) rectangle ++(3,-3);
        \fi
      }}

    \foreach \r [count=\i from 0] in \griglia {
      \foreach \v [count=\j from 0] in \r {
        \node[inchiostro] at (\j+0.5,-\i-0.5) {\v};
      }}

    \foreach \k in {0,...,9}{
      \draw[grigiotratto,line width=0.5pt] (\k,0) -- (\k,-9);
      \draw[grigiotratto,line width=0.5pt] (0,-\k) -- (9,-\k);
    }
    \foreach \k in {0,3,6,9}{
      \draw[blunotte,line width=1.1pt] (\k,0) -- (\k,-9);
      \draw[blunotte,line width=1.1pt] (0,-\k) -- (9,-\k);
    }
    \node[grigiotesto,font=\footnotesize\sffamily,anchor=north] at (4.5,-9.3)
      {la griglia};
  \end{scope}

  % ---- i determinanti dei blocchi -----------------------------------
  \begin{scope}[xshift=11.2cm]
    % celle larghe 2 e alte 2.5, per allinearsi ai blocchi della griglia
    \foreach \d [count=\n from 0] in {0,-177,81,9,11,-3,-243,-5,77}{
      \pgfmathtruncatemacro{\bi}{div(\n,3)}
      \pgfmathtruncatemacro{\bj}{mod(\n,3)}
      \node[rossoterra] at (2*\bj+1,{-0.75-2.5*\bi}) {$\d$};
    }
    \foreach \k in {2,4}{
      \draw[grigiotratto,line width=0.5pt] (\k,0.5) -- (\k,-7);
    }
    \foreach \k in {1,2}{
      \draw[grigiotratto,line width=0.5pt] (0,{0.5-2.5*\k}) -- (6,{0.5-2.5*\k});
    }
    \draw[blunotte,line width=1.1pt] (0,0.5) rectangle (6,-7);
    \node[grigiotesto,font=\footnotesize\sffamily,anchor=north,align=center]
      at (3,-7.4)
      {determinante di ciascun blocco\\[2pt] somma: $-250$};
  \end{scope}

  % ---- didascalia complessiva ---------------------------------------
  \node[grigiotesto,font=\footnotesize\sffamily,anchor=north] at (8.6,-10.6)
    {determinante della griglia intera: $-2\,975\,130=-2\cdot3^4\cdot5\cdot3673$};
\end{tikzpicture}

  \caption{La griglia del quaderno. Accanto, in rosso, il determinante di ciascun blocco $3\times3$: quello in alto a sinistra è nullo, cioè il blocco è singolare.}
  \label{fig:sudoku_griglia}
\end{figure}

I nove valori sono $0$, $-177$, $81$, $9$, $11$, $-3$, $-243$, $-5$, $77$, e la loro somma è $-250$. Tutti i valori sono stati ricalcolati durante il recupero e coincidono, compreso lo zero del primo blocco — che è l'unico caso interessante del gruppo, perché segnala che le tre righe di quel blocco sono linearmente dipendenti.

La somma $-250$ non ha alcun rapporto con il determinante della griglia intera, e non c'è motivo perché ne abbia: il determinante di una matrice a blocchi non è una funzione dei determinanti dei blocchi, a meno che la matrice non sia triangolare a blocchi. Il quaderno lo verifica sperimentalmente, che è un modo legittimo di rispondersi.

\section{Le nove righe rimodellate}

Seconda idea, diversa dalla prima: invece dei blocchi, si prende ogni riga della griglia e la si ripiega in una matrice $3\times3$. Nove righe, nove matrici, nove determinanti.

\begin{espr}{23}{User}{}
\[
\begin{vmatrix} 9 & 1 & 3 \\ 7 & 8 & 6 \\ 5 & 4 & 2 \end{vmatrix}
\]
\end{espr}
La prima riga della griglia ripiegata in una matrice $3\times3$.

\begin{espr}{24}{Simp(\#23)}{}
\[
-92
\]
\end{espr}
$-92$.

\begin{espr}{25}{User}{}
\[
\begin{vmatrix} 4 & 8 & 7 \\ 5 & 2 & 3 \\ 1 & 9 & 6 \end{vmatrix}
\]
\end{espr}
La seconda riga.

\begin{espr}{26}{Simp(\#25)}{}
\[
25
\]
\end{espr}
$25$.

\begin{espr}{27}{User}{}
\[
\begin{vmatrix} 2 & 6 & 5 \\ 4 & 1 & 9 \\ 3 & 8 & 7 \end{vmatrix}
\]
\end{espr}
La terza riga. Il valore, $9$, è quello usato nella somma della riga $\#41$.

\begin{espr}{28}{Simp(\#27)}{}
\[
9
\]
\end{espr}
\begin{espr}{29}{User}{}
\[
\begin{vmatrix} 8 & 5 & 6 \\ 3 & 7 & 4 \\ 2 & 1 & 9 \end{vmatrix}
\]
\end{espr}
La quarta riga.

\begin{espr}{30}{Simp(\#29)}{}
\[
311
\]
\end{espr}
$311$.

\begin{espr}{31}{User}{}
\[
\begin{vmatrix} 7 & 9 & 2 \\ 1 & 6 & 5 \\ 4 & 3 & 8 \end{vmatrix}
\]
\end{espr}
La quinta riga.

\begin{espr}{32}{Simp(\#31)}{}
\[
297
\]
\end{espr}
$297$.

\begin{espr}{33}{User}{}
\[
\begin{vmatrix} 3 & 4 & 1 \\ 2 & 9 & 8 \\ 7 & 6 & 5 \end{vmatrix}
\]
\end{espr}
La sesta riga.

\begin{espr}{34}{Simp(\#33)}{}
\[
124
\]
\end{espr}
$124$.

\begin{espr}{35}{User}{}
\[
\begin{vmatrix} 6 & 3 & 4 \\ 9 & 5 & 7 \\ 8 & 2 & 1 \end{vmatrix}
\]
\end{espr}
La settima riga.

\begin{espr}{36}{Simp(\#35)}{}
\[
-1
\]
\end{espr}
$-1$.

\begin{espr}{37}{User}{}
\[
\begin{vmatrix} 5 & 2 & 9 \\ 8 & 4 & 1 \\ 6 & 7 & 3 \end{vmatrix}
\]
\end{espr}
L'ottava riga.

\begin{espr}{38}{Simp(\#37)}{}
\[
277
\]
\end{espr}
$277$.

\begin{espr}{39}{User}{}
\[
\begin{vmatrix} 1 & 7 & 8 \\ 6 & 3 & 2 \\ 9 & 5 & 4 \end{vmatrix}
\]
\end{espr}
La nona riga.

\begin{espr}{40}{Simp(\#39)}{}
\[
-16
\]
\end{espr}
$-16$.

\begin{espr}{41}{User}{}
\[
-92 + 25 + 9 + 311 + 297 + 124 - 1 + 277 - 16
\]
\end{espr}
La somma dei nove valori.

\begin{espr}{42}{Simp(\#41)}{}
\[
934
\]
\end{espr}
Vale $934$.


I nove valori sono $-92$, $25$, $9$, $311$, $297$, $124$, $-1$, $277$, $-16$, e sommano a $934$. Anche questi coincidono tutti con il ricalcolo.

Né $-250$ né $934$ hanno un rapporto riconoscibile con $-2\,975\,130$, e nessuno dei due è multiplo di $45$: la proprietà di divisibilità vale per la griglia intera perché dipende dalle somme di righe e colonne, che nei sottoblocchi non sono $45$.

\section*{In sintesi}
\addcontentsline{toc}{section}{In sintesi}

Il quaderno pone una domanda e ottiene una risposta negativa: i determinanti dei sottoblocchi di un sudoku non dicono nulla sul determinante della griglia, e le loro somme non sono invarianti significativi. È una risposta utile quanto una affermativa, e per ottenerla servivano diciannove determinanti calcolati a mano — che nel 2009, con Derive aperto sullo schermo, erano diciannove righe.

